\chapter{Examples}
\section{Citation and Referencing}
Referece: Appendix~\ref{appendix-a:title}

\subsubsection{\textbackslash cite}
\cite{ombudsman2024machinery} \\

\subsubsection{\textbackslash citep}
\citep{ombudsman2024machinery}

\subsubsection{\textbackslash citet}
\citet{dang2023toward}





\import{./sections}{section-example}

\section{Macro Help}
\import{./partials/}{macro-help}




%====
\section{Example Table}
Note: Instead of \textbackslash{caption}, use \textbackslash{utsTableCaption} - this allows the caption to be left aligned\\
\\
I recommend following this pattern. Put your table in a "partial" file, then reference it like this:
\import{./partials/}{table-example}




%====
\section{Example Table with paragraph text}
Note: Instead of \textbackslash{caption}, use \textbackslash{utsTableCaption} - this allows the caption to be left aligned\\
\\
I recommend following this pattern. Put your table in a "partial" file, then reference it like this:
\import{./partials/}{table-paragraph-text-example.tex}




%====
\section{Example Long Table spilling over to another page}
Note: use \textbackslash{caption} as normal, but use the \textbackslash{utsLongTableCaptionSetup} macro before the longtable - this allows the caption to be left aligned\\
\\
I recommend following this pattern. Put your table in a "partial" file, then reference it like this:
\import{./partials/}{table-long-example.tex}




%====
\section{Example Figure}
Special note: Avoid using PNG, JPG, etc. It is better to use PDF to preserve resolution.

I recommend following this pattern. Put your figure in a "partial" file, then reference it like this:
\import{./partials/}{figure-example}



\import{./partials/}{indexing}


